\documentclass[11pt,a4paper]{article}
\usepackage[margin=25mm]{geometry}
\usepackage{amsmath,amssymb,bm}
\usepackage{booktabs}
\usepackage{array}
\usepackage[hidelinks]{hyperref}

\newcommand{\bsig}{\bm{\sigma}}
\newcommand{\beps}{\bm{\varepsilon}}
\newcommand{\bfeta}{\bm{\eta}}
\newcommand{\bG}{\bm{G}}
\newcommand{\ihat}{\hat{\bm{\imath}}}
\newcommand{\bn}{\bm{n}}
\newcommand{\tr}{\operatorname{tr}}
\newcommand{\dev}{\operatorname{dev}}

\title{A reduced return mapping for the cohesive visco-elastic model}
\author{Technical note (FEMCubes, \texttt{Solid:Cohesive\_ViscoElastic})}
\date{November 2026}

\begin{document}
\maketitle

\section{Purpose}

The cohesive visco-elastic model combines elasticity, power-law creep and fracture eigenstrains.
At every integration point, the stress, the creep multiplier and the fracture multipliers follow
from a small set of nonlinear equations, the local problem. The C++ kernel solves this problem
in full tensor form (six stress components plus the multipliers) by a Newton method. In fracture
simulations this local solve takes around 75--80\,\% of the run time.

This note shows that the local problem only has two or three scalar unknowns. The material is
isotropic, and all strain directions in the model are built from the volumetric and deviatoric
parts of the strain. As a result, the stress, the creep strain and the fracture eigenstrain all
lie in a plane spanned by two fixed directions. Most integration points then need no iteration
at all, and the remaining points solve a system of two or three equations. This gives the same
answer as the kernel, but at a cost of about $0.25$--$0.5\,\mu$s per point and call instead of
$12$--$15\,\mu$s. The reference implementation is \texttt{src/femcubes/solids/\_reduced\_cve.py}.

\section{The local problem}

We use the notation of the model paper. Within a time increment $\Delta t$, the total strain
$\beps$ and the creep strain of the previous step $\beps^v_\mathrm{old}$ are known, as is the
degradation $d=d(\phi)\in[\kappa,1]$. The elastic strain is what remains after subtracting the
creep strain and the fracture eigenstrain,
\begin{equation}
  \bsig = \mathbb{D}:\left(\beps-\beps^v_\mathrm{old}-\Delta\lambda_v\,\frac{\partial f_v}{\partial\bsig}-\bfeta\right),
  \label{eq:stress}
\end{equation}
with $\mathbb{D}$ the isotropic elastic stiffness. The creep strain increment follows the
power-law flow rule, written for the creep multiplier as
\begin{equation}
  \Delta\lambda_v = \Delta t\left(\frac{f_v(\bsig)}{\sigma_c}\right)^n,\qquad
  f_v(\bsig)=\sqrt{\tfrac{n_d}{2}\,\bsig_\mathrm{dev}:\bsig_\mathrm{dev}},\qquad
  \sigma_c=\left(\frac{n_d^{(n+1)/2}}{2A}\right)^{1/n},
  \label{eq:creep}
\end{equation}
where $n_d$ is the number of dimensions ($n_d=2$ for the intrinsic two-dimensional model,
$n_d=3$ otherwise) and $A$ and $n$ are the parameters of Glen's law. The fracture eigenstrain
is built from one or two multipliers and their direction tensors,
\begin{equation}
  \bfeta=\sum_i\lambda_i\,\bG_i,\qquad \lambda_i\geq0 .
\end{equation}
For each multiplier, either the multiplier is zero, or the stress lies on the strength surface
in its direction. With the drive
\begin{equation}
  r_i = \bG_i:\left(\frac{\partial F}{\partial\bfeta}-\bsig\right)+\zeta\,\lambda_i ,
  \label{eq:drive}
\end{equation}
this reads $\lambda_i\geq0$, $r_i\geq0$ and $\lambda_i\,r_i=0$. Here $F$ is the total fracture
potential (strength and interpenetration parts, scaled by $d$ where applicable) and $\zeta$ a
small cohesive viscosity that the kernel adds for robustness ($\zeta=10^{-8}K\kappa$).

The direction tensors are built from the strain that remains after creep,
$\bm{s}=\beps-\beps^v_\mathrm{old}-\Delta\lambda_v\,\partial f_v/\partial\bsig$:
\begin{itemize}
  \item failure types with two multipliers (\texttt{r1}, \texttt{r1\_DP},
    \texttt{r2\_2L}, \texttt{r2\_2L\_DP}) use $\bG_1=\pm\bm{I}/\sqrt{n_d}$ (sign of
    $\tr\bm{s}$) and $\bG_2=\dev\bm{s}/\|\dev\bm{s}\|$;
  \item failure types with one multiplier (\texttt{r2}, \texttt{r2\_DP}) use $\bG=\bm{s}/\|\bm{s}\|$ in tension and
    $\bG=\dev\bm{s}/\|\dev\bm{s}\|$ in compression.
\end{itemize}

\section{Reduction to two strain invariants}

We split the trial strain $\beps-\beps^v_\mathrm{old}$ into its volumetric and deviatoric parts,
\begin{equation}
  \beps-\beps^v_\mathrm{old}=e_v\,\ihat+e_q\,\bn,\qquad
  \ihat=\frac{\bm{I}}{\sqrt{n_d}},\qquad
  \bn=\frac{\dev(\beps-\beps^v_\mathrm{old})}{\|\dev(\beps-\beps^v_\mathrm{old})\|},
\end{equation}
so that $e_v=\ihat:(\beps-\beps^v_\mathrm{old})$ and $e_q\geq0$. Both $\ihat$ and $\bn$ have unit
norm and are orthogonal. Suppose now that the stress lies in their plane,
\begin{equation}
  \bsig=a\,\ihat+b\,\bn .
\end{equation}
Then the creep direction is $\partial f_v/\partial\bsig=\gamma\,\bn$ (for $b>0$), with
$\gamma=\sqrt{n_d/2}$, and $f_v=\gamma\,|b|$. The strain after creep becomes
$\bm{s}=e_v\,\ihat+(e_q-\gamma\Delta\lambda_v)\,\bn$, which lies in the same plane. All direction
tensors above are built from $\bm{s}$, so they also lie in this plane, and so does the
eigenstrain,
\begin{equation}
  \bfeta=\eta_v\,\ihat+\eta_q\,\bn .
\end{equation}
The elastic stiffness maps $\ihat$ onto $k_v\,\ihat$ and $\bn$ onto $2\mu\,\bn$, with
$k_v=n_dK$ ($K$ the $n_d$-dimensional bulk modulus, $k_v=n_d\lambda+2\mu$). Eq.~\eqref{eq:stress}
therefore indeed gives a stress in the same plane:
\begin{equation}
  a=k_v\,(e_v-\eta_v),\qquad b=2\mu\,(e_q-\gamma\,\Delta\lambda_v-\eta_q).
  \label{eq:ab}
\end{equation}
The assumed form is thus consistent, and since the local problem has a unique solution for the
standard (associated) failure types, it is the solution. The full problem therefore reduces to
the scalar unknowns $\Delta\lambda_v$ and $\lambda_1$ (and $\lambda_2$), with $(e_v,e_q)$, $d$
and $\Delta t$ as the only inputs. The creep update becomes
$\Delta\lambda_v=\Delta t\,(\gamma b/\sigma_c)^n$.

In these coordinates, the trace and deviatoric norm of the eigenstrain are
$\tr\bfeta=\sqrt{n_d}\,\eta_v$ and $\|\dev\bfeta\|=|\eta_q|$, and the projections of the stress
on the directions are $\bsig:\bG_i=a\,G_{i,v}+b\,G_{i,q}$. The directions themselves become
\begin{align}
  &\text{two multipliers:} && \bG_1=(\pm1,\,0),\qquad \bG_2=(0,\,1),\\
  &\text{one multiplier, tension:} && \bG=\frac{(e_v,\;e_q-\gamma\Delta\lambda_v)}{\sqrt{e_v^2+(e_q-\gamma\Delta\lambda_v)^2}},
\end{align}
written as (volumetric, deviatoric) components.
It should be noted that the single-multiplier direction in tension depends on the creep
multiplier. This couples the creep and fracture equations, which is why these types need the
general solver of Section~\ref{sec:return} and not the closed-form route of
Section~\ref{sec:r2}.

\section{Fracture potentials in reduced form}

Every failure potential of the model is a function of $\tr\bfeta$, $\|\dev\bfeta\|$ and, through
its branches, of the trace of the strain after creep, $\tr\bm{s}=\sqrt{n_d}\,e_v$. In reduced
form, the potentials are therefore functions $F(\eta_v,\eta_q;e_v,d)$. With
\begin{equation}
  F_1=\frac{f}{n_d}\tr\bfeta=\frac{f}{\sqrt{n_d}}\,\eta_v,\qquad
  F_2=\frac{f_s}{\sqrt2}\,\|\dev\bfeta\|,
\end{equation}
where $f$ is $f_t$ or the calibrated $f^*$ of the split r2 criteria, Table~\ref{tab:F} lists the
potentials as implemented in the kernel (\texttt{shared/fracture\_relations.h}). In compression
($e_v<0$) a large penalty on $\tr\bfeta$ prevents interpenetration. For the Drucker--Prager
variants, the deviatoric strength $f_s$ is replaced by a confined strength
$q=\min(f_s-K\tan\phi\,\tr\bm{s},\,q_\mathrm{max})$ (\texttt{r2\_2L\_DP} uses the pressure
$\max(-K\tr(\bm{s}-\bfeta),0)$ instead). The kernel regularises
$\|\dev\bfeta\|$ as $\sqrt{\eta_q^2+10^{-20}}-10^{-10}$ and smooths the r2 norm with
$\varepsilon=10^{-6}$; the reduced form keeps both, so the two implementations solve the same
equations.

\begin{table}[ht]
  \centering
  \small
  \begin{tabular}{>{\ttfamily}l l l}
    \toprule
    \textrm{FailureType} & tension ($e_v\geq0$) & compression ($e_v<0$) \\
    \midrule
    r1 & $d\,(F_1+F_2)$ & $-10^6F_1+d\,F_2$ \\
    r2, r2\_2L, r2\_2L\_DP & $d\sqrt{F_1^{*2}+F_2^{*2}}$ & $-10^6F_1^*+d\,F_2^*$ \\
    r2\_DP & $d\sqrt{F_1^2+F_2^2}$ & $-10^{12}F_1+d\,F_2\,q/f_s$ \\
    r1\_DP & $d\,F_1+d\,F_2\,q/f_s$ & $-10^6F_1+d\,F_2\,q/f_s$ \\
    \bottomrule
  \end{tabular}
  \caption{Fracture potentials $F$ in reduced form. $F^*$ uses the r2 calibration of the model
  paper ($F_1^*=f_t^*\tr\bfeta$, $F_2^*=f_s^*\|\dev\bfeta\|$; the split types calibrate $f_t^*$
  against uniaxial tension). For \texttt{r2\_2L\_DP}, $f_s^*$ is confined in compression; for
  \texttt{r1\_DP}, $q$ takes its unconfined value $q_0=\max(f_s-f_t\tan\phi/3,0)$ when
  $K\tr\bm{s}\geq f_t/3$.}
  \label{tab:F}
\end{table}

With the gradient $F_{,v}=\partial F/\partial\eta_v$, $F_{,q}=\partial F/\partial\eta_q$, the
drive \eqref{eq:drive} becomes
\begin{equation}
  r_i=G_{i,v}\,(F_{,v}-a)+G_{i,q}\,(F_{,q}-b)+\zeta\,\lambda_i .
  \label{eq:drive_reduced}
\end{equation}
This is the full content of the fracture model in the reduced coordinates. Any other isotropic
strength surface can be added by providing $F$, its first and second derivatives with respect to
$(\eta_v,\eta_q)$, and the derivative of its gradient with respect to $e_v$ (for
strain-dependent strengths).

\section{The creep-only state}

Most integration points do not crack. For these, $\bfeta=0$ and only creep remains. The
deviatoric stress is then $b=2\mu\,(e_q-\gamma\Delta\lambda_v)$, and with
$x=\gamma\Delta\lambda_v/e_q$ the creep update \eqref{eq:creep} becomes the scalar equation
\begin{equation}
  x=c\,(1-x)^n,\qquad
  c=\frac{\gamma\,\Delta t}{e_q}\left(\frac{2\mu\,\gamma\,e_q}{\sigma_c}\right)^n .
  \label{eq:creep_only}
\end{equation}
The right-hand side decreases with $x$ and the left-hand side increases, so there is exactly one
root in $[0,1)$. This root depends on the single number $c$, the creep strain the trial stress
would cause in one step, relative to $e_q$. We solve it once per creep exponent, by bisection to
round-off, and tabulate $s=-\ln(1-x)$ against $\log_{10}c$ (cubic Hermite interpolation with
exact derivatives, $10^{-30}\leq c\leq10^{14}$, $10^{-3}$ decades). Outside the table the
asymptotes $s\approx c$ and $s\approx(\ln c)/n$ are used. The creep-only stress is then
\begin{equation}
  a^0=k_v\,e_v,\qquad b^0=2\mu\,e_q\,e^{-s},\qquad \Delta\lambda_v^0=\frac{e_q\,(1-e^{-s})}{\gamma},
\end{equation}
and its derivative with respect to $e_q$ follows from the tabulated $\mathrm{d}s/\mathrm{d}\log_{10}c$
and $\mathrm{d}\log_{10}c/\mathrm{d}e_q=(n-1)/(e_q\ln10)$. No iteration is needed.

\section{Which points crack}

The fracture modes are tested at the creep-only state, as in the kernel. A mode is active when its
drive \eqref{eq:drive_reduced}, evaluated at a small probe multiplier ($10^{-9}$), is negative
(below $-10^{-6}\,d\max(f_t,f_s)$).

For the split r2 criteria in tension, this probe is not a valid test. The admissible stresses of
$F=d\sqrt{F_1^{*2}+F_2^{*2}}$ form an ellipse,
\begin{equation}
  \left(\frac{a}{d\,c_1}\right)^2+\left(\frac{b}{d\,c_2}\right)^2\leq1,\qquad
  c_1=\sqrt{n_d}\,f_t^*,\qquad c_2=f_s^*,
  \label{eq:ellipse}
\end{equation}
and a point cracks exactly when its creep-only stress lies outside this ellipse. Testing each mode
separately misses points that are within both single-mode strengths but outside the ellipse.
The probe drives, in turn, see the gradient of the smoothed norm along the probe direction and
flag points inside the ellipse; the kernel then converges there to tiny ``regularised''
multipliers of order $10^{-7}f_t/E$, which the sharp potential does not have. The reduced update
therefore uses the ellipse test \eqref{eq:ellipse} alone in tension, and returns exactly zero
multipliers inside the surface.

Points without an active mode keep the creep-only state. For the associated failure types the
local problem minimises a strictly convex incremental potential, so this test is exact.

\section{The reduced return}
\label{sec:return}

Points with an active mode solve the reduced system for $\bm{u}=(\Delta\lambda_v,\lambda_1[,\lambda_2])$:
\begin{equation}
  R_0=\Delta\lambda_v-\Delta t\left(\frac{\gamma\langle b\rangle}{\sigma_c}\right)^n=0,\qquad
  R_i=\varphi\!\left(2\mu\,\lambda_i,\;r_i\right)=0,
  \label{eq:system}
\end{equation}
where $a$ and $b$ follow from Eq.~\eqref{eq:ab} and $\varphi(x,y)=\sqrt{x^2+y^2}-x-y$ is the
Fischer--Burmeister function. It is zero exactly when $x\geq0$, $y\geq0$ and $x\,y=0$, so each
$R_i=0$ states that either the multiplier is zero or the drive vanishes. The multiplier is scaled
by $2\mu$ to give both arguments units of stress, as in the kernel.

The system is solved by a semi-smooth Newton method with a backtracking line search on
$\tfrac12\|(2\mu R_0,R_1,R_2)\|^2$, keeping the multipliers non-negative. It starts from the
creep-only state, with the active multipliers estimated from one elastic step,
$\lambda_i=-r_i^0/(k_vG_{i,v}^2+2\mu G_{i,q}^2)$. The Jacobian is available in closed form.
With $\bm{H}$ the Hessian of $F$ with respect to $(\eta_v,\eta_q)$,
$\bm{k}=\mathrm{diag}(k_v,2\mu)$ and $\bm{w}=(F_{,v}-a,\,F_{,q}-b)$,
\begin{align}
  \frac{\partial r_i}{\partial\lambda_j}&=\bG_i\cdot(\bm{H}+\bm{k})\,\bG_j+\zeta\,\delta_{ij},\\
  \frac{\partial r_i}{\partial\Delta\lambda_v}&=\bG_i\cdot\left[(\bm{H}+\bm{k})\,\frac{\partial\bfeta}{\partial\Delta\lambda_v}+(0,\,2\mu\gamma)\right]+\bm{w}\cdot\frac{\partial\bG_i}{\partial\Delta\lambda_v},
\end{align}
where $\partial\bfeta/\partial\Delta\lambda_v=\sum_j\lambda_j\,\partial\bG_j/\partial\Delta\lambda_v$
is non-zero only for the single-multiplier directions in tension. The creep row follows from
$\partial b/\partial\lambda_j=-2\mu\,G_{j,q}$ and
$\partial b/\partial\Delta\lambda_v=-2\mu\,(\gamma+\partial\eta_q/\partial\Delta\lambda_v)$.
The system has two or three unknowns and is solved by Gaussian elimination. Newton converges
quadratically, typically in three to six iterations, to a scaled residual of $10^{-12}$ of the
stress scale.

If the method does not converge (rare, and so far only seen in extreme states, with strains of
tens of times the strength strain in nearly fully damaged material), the point is recomputed by
the C++ kernel with the inputs the kernel would otherwise have received.

\section{Consistent tangent}

The stress depends on the strain through $(e_v,e_q)$ and through the direction $\bn$. The scalar
part follows by implicit differentiation of the converged system \eqref{eq:system}: with
$\bm{J}_u=\partial\bm{R}/\partial\bm{u}$ and $\bm{J}_e=\partial\bm{R}/\partial(e_v,e_q)$ at fixed
$\bm{u}$,
\begin{equation}
  \frac{\mathrm{d}\bm{u}}{\mathrm{d}(e_v,e_q)}=-\bm{J}_u^{-1}\bm{J}_e,\qquad
  \frac{\mathrm{d}a}{\mathrm{d}e_\alpha}=\frac{\partial a}{\partial e_\alpha}+\frac{\partial a}{\partial\bm{u}}\cdot\frac{\mathrm{d}\bm{u}}{\mathrm{d}e_\alpha},
\end{equation}
and likewise for $b$. At the converged state, the Fischer--Burmeister rows reduce to
$\mathrm{d}r_i=0$ for active modes and $\mathrm{d}\lambda_i=0$ for inactive ones. Since
$\mathrm{d}e_v/\mathrm{d}\beps=\ihat$, $\mathrm{d}e_q/\mathrm{d}\beps=\bn$ and
$\mathrm{d}\bn/\mathrm{d}\beps=(\mathbb{P}_\mathrm{dev}-\bn\otimes\bn)/e_q$, the tangent is
\begin{equation}
  \frac{\mathrm{d}\bsig}{\mathrm{d}\beps}=
  \ihat\otimes\left(\frac{\mathrm{d}a}{\mathrm{d}e_v}\ihat+\frac{\mathrm{d}a}{\mathrm{d}e_q}\bn\right)
  +\bn\otimes\left(\frac{\mathrm{d}b}{\mathrm{d}e_v}\ihat+\frac{\mathrm{d}b}{\mathrm{d}e_q}\bn\right)
  +\frac{b}{e_q}\left(\mathbb{P}_\mathrm{dev}-\bn\otimes\bn\right).
\end{equation}
The last term describes the rotation of the deviatoric direction with the strain; it is the
elastic shear stiffness for creep-only points without creep. The tangent is the exact
derivative of the update; it agrees with finite differences to about $10^{-7}$ of the elastic
stiffness for every failure type.

\section{Plane stress}

In plane stress the out-of-plane strain $e_{zz}$ is an additional unknown, set by
$\sigma_{zz}=0$. Since $\sigma_{zz}$ increases with $e_{zz}$, a safeguarded scalar Newton method
is used per point: Newton steps with $\mathrm{d}\sigma_{zz}/\mathrm{d}e_{zz}$ from the tangent,
inside a bracket of known signs, with bisection when a step leaves it. It starts from the elastic
plane-stress strain, so elastic points need no iteration. Each evaluation is the reduced
three-dimensional update above. The solved $e_{zz}$ is returned, and the tangent is condensed as
in the kernel,
\begin{equation}
  \tilde{\mathbb{C}}=\mathbb{C}-\frac{(\mathbb{C}:\bm{e}_z\otimes\bm{e}_z)\otimes(\bm{e}_z\otimes\bm{e}_z:\mathbb{C})}{C_{zzzz}},
\end{equation}
with zero $zz$ row and column.

\section{The split r2 criterion as a single angle}
\label{sec:r2}

For \texttt{r2\_2L} the return has an even simpler structure. When a point returns in tension,
the stress lands on the ellipse
\eqref{eq:ellipse}, so it can be written with one angle,
\begin{equation}
  (a,b)=d\,(c_1\cos\theta,\;c_2\sin\theta)
\end{equation}
(corrected for the small cohesive viscosity). The creep follows from $b$, and the angle is fixed
by normality: the return vector, from the trial stress (with the creep of the returned stress)
to $(a,b)$, must be parallel to the outward normal of the ellipse in the energy metric,
\begin{equation}
  \bm{v}(\theta)=\left(\frac{k_v\cos\theta}{c_1},\;\frac{2\mu\sin\theta}{c_2}\right).
\end{equation}
Their cross product $R(\theta)$ is negative at
$\theta=0$ and positive at $\theta=\pi/2$, so a bracketed scalar Newton method finds the angle in
about five iterations. The general solver of Section~\ref{sec:return} gives the same result.

\section{Accuracy and cost}

The reduced update was compared with the C++ kernel at random states (strains from $0.1$ to
$30$ times the strength strain, damage down to $10^{-3}$) for every failure type, in intrinsic
two dimensions, three dimensions and plane stress. The stress agrees to the kernel's tolerance.
Where the two differ, at strongly damaged points with large openings, the kernel's accepted
states are the less accurate ones: they violate the creep equation \eqref{eq:creep}, or, in plane
stress, the kernel's own three-dimensional response at the same strain. The reduced update
satisfies both to round-off.

In coupled simulations (fine plate with hole, \texttt{r2\_2L}), the reduced update reproduces the
kernel's load--displacement response to about $10^{-7}$ of the peak force, with the same time
steps, failed steps and crack. It costs $0.24\,\mu$s per point and call in two dimensions and
$0.47\,\mu$s in plane stress, against $11.5$--$15\,\mu$s for the kernel. The run times drop by
factors of $5.6$ (intrinsic two dimensions), $6.4$ (plane strain) and $6.7$ (plane stress); the
remaining time is mostly the linear solves. On a coarser plate, the other failure types that
crack in this test (\texttt{r1}, \texttt{r1\_DP}, \texttt{r2}, \texttt{r2\_DP},
\texttt{r2\_2L\_DP}) also reproduce the kernel, with the same time steps, failed steps and crack
length and a load--displacement response within $3\cdot10^{-7}$ of the peak force.

\end{document}
